\documentclass[11pt]{beamer}

\usetheme{m}

\usepackage{booktabs}
\usepackage[scale=2]{ccicons}

% \usepackage{pgfplots}
% \usepgfplotslibrary{dateplot}

\usepackage[latin1]{inputenc}
\usepackage[spanish]{babel}

\usepackage{listings}
\usepackage{pgfplots}

\usepackage{amsmath}

\title{Maui}
\subtitle{Cluster Scheduler}
\date{}
\author{Luis Mar�a Costero Valero\\Jes�s Javier Dom�nech Arellano\\Hristo
  Ivanov Ivanov}
\institute{26 Enero 2016}

% \titlegraphic{\hfill\includegraphics[scale=0.2]{HPFF_logo.jpg}}


%----------------------------------------------------------------------
% Modificaciones del template
\metroset{titleformat=regular}
\metroset{inner/sectiontitleformat=regular}
\metroset{outer/frametitleformat=regular}
\metroset{block=fill}
\newcommand\sectionDark[1]{{\metroset{background=dark} \section{#1} }}
\newcommand\sectionDarkSpecial[1]{{\metroset{background=dark} \section*{#1}
  }}





\begin{document}

\maketitle

%% ======= INDICE ========================================
% \begin{frame}
%   \frametitle{�ndice}
%   \setbeamertemplate{section in toc}[sections numbered]
%   \tableofcontents[hideallsubsections]
% \end{frame}

% \section{Introduction} %=== Esrto genera una p�gina de secci�n.

\begin{frame}
  \frametitle{�Qu� es?}
  
  \begin{block}{Maui Scheduler}
    ~\\
    -- Planificador de tareas para clusters.\\
    -- Desarrollo a mediados de los 90.\\
    -- Desarrollado por \texttt{Cluster Resources, Inc.}\\
    -- C�digo libre. Licencia no-comercial.\\
    -- Actualmente en desuso.\\
    Recomiendan migrar a \emph{Moab} (c�digo cerrado).\\
    ~
  \end{block}
\end{frame}


\begin{frame}
  \frametitle{Caracter�sticas de Maui I}
  \begin{enumerate}
  \item M�ltiples pol�ticas de planificaci�n.
  \item Ajuste de prioridades de forma din�mica.
  \item Reserva de nodos.
  \item Sistema de control de procesos avanzado.
  \end{enumerate}
\end{frame}


\begin{frame}
  \frametitle{Caracter�sticas de Maui II}
  \begin{enumerate}
  \item \textbf{Control de Tr�fico:} Gestiona las tareas para que no
    interfieran entre si. 
  \item \textbf{Centrado en Objetivos:} Configuraci�n de cluster seg�n la tarea.
  \item \textbf{Optimizaciones:} Supeditado a los dos puntos
    anteriores, el planificador intenta aprovechar los recursos del cluster.
  \end{enumerate}
\end{frame}

\begin{frame}
  \frametitle{Job}
  
  \textbf{Job: } La unidad de trabajo de Maui (lo que va a ejecutar).\\
  Un \emph{job} est� formado a su vez por:

  \begin{enumerate}
  \item Due�o del trabajo.
  \item Varios \textbf{Requisitos} (cada uno con sus limitaciones de recursos).
  \item Tiempo m�ximo de ejecuci�n.
  \end{enumerate}
 
\end{frame}


\begin{frame}
  \frametitle{Requisitos}
  Un \textbf{Requisito} est� formado por:
  \begin{itemize}
  \item Definici�n de la tarea.
  \item L�mite de recursos: Procesadores, memoria, swap, disk.
  \item Task count: N� de instancias de la tarea a ejecutar en el requisito.
  \item Task list: Lista de nodos en los que la tarea tiene que ser
    ejecutada.
  \item Recursos estad�sticos: Qu� aspectos de la tarea monitorear.
  \end{itemize}

  \begin{block}{Ejemplo}
    Un \textbf{trabajo} puede consistir en 2 requisitos:
    \begin{enumerate}
    \item 1 nodo, 512 MB RAM como m�nimo, ejecutar la tarea X
    \item 24 nodos, 128 MB como m�nimo, ejecutar la tarea Y
    \end{enumerate}
  \end{block}
\end{frame}

\begin{frame}
  \frametitle{Tipos de Nodos}
  Maui contempla los siguientes tipos de nodos:
  \begin{enumerate}
  \item \textbf{SHARED: } Se puede ejecutar cualquier combinaci�n de tareas
    provenientes de cualquier job.
  \item \textbf{SINGLEUSER: } Se pueden ejecutar distintas tareas siempre
    que sean del mismo usuario.
  \item \textbf{SINGLEJOB: } Se pueden ejecutar distintas tareas siempre
    que sean del mismo job.
  \item \textbf{SINGLETASK: } Solamente se puede ejecutar una tarea simult�neamente.
  \end{enumerate}
\end{frame}



\begin{frame}
  \frametitle{Prioridad a las tareas:}
  La prioridad se calcula de acuerdo a cada \alert{componente}:\\
  \begin{center}
    \texttt{Nodos}, \texttt{Procesadores}, \texttt{Memoria}, \texttt{Swap}, \texttt{Disco}
  \end{center}

  \begin{center}
    \texttt{<COMPONENT WEIGHT> * <SUBCOMPONENT WEIGHT> * <PROPRITY
      SUBCOMPONENT VALUE>}
  \end{center}

  \begin{enumerate}
  \item Gesti�n de trabajos mediante una FIFO de prioridades.
  \item Posibilidad de modificar las prioridades en tiempo real.
  \end{enumerate}
\end{frame}


\begin{frame}
  \frametitle{Asignaci�n de nodos}
  Maui proporciona distintas pol�ticas para la asignaci�n de nodos:

  \begin{itemize}
  \item \textbf{CPULOAD: } M�ximo n�mero de cpus sin usar.
  \item \textbf{FIRSTAVAILABLE: } Primer nodo libre.
  \item \textbf{MINRESOURCE: } Nodo que desaproveche menos recursos.
  \item \textbf{CONTIGUOUS: } Asignar nodos contiguos.
  \item \textbf{MAXBALANCE: } Intenta buscar los mejores nodos para el
    trabajo (lento).
  \item \textbf{FASTEST: } Selecci�n de los nodos m�s r�pidos.
  \item \textbf{LOCAL: } Algoritmo propio de cada nodo (opcional).
  \end{itemize}
\end{frame}



\begin{frame}
  \frametitle{Bibliograf�a}
  \begin{enumerate}
  \item ``Maui Scheduler Administrator's Guide'' \\\url{www.clusterresources.com/products/maui/docs/}
  \item ``Maui Scheduler'' \\\url{www.clusterresources.com/products/maui/docs/}
  \item Art�culo Wikipedia\\\url{en.wikipedia.org/wiki/Maui_Cluster_Scheduler}
  \end{enumerate}
\end{frame}
\end{document}

% 
% 
%%% 
%%% Local Variables:
%%% mode: latex
%%% TeX-master: "./PRESENTACION.tex"
%%% End:
